\documentclass[a4paper, noanswers, addpoints, 12pt]{exam}
\usepackage[T1]{fontenc}
\usepackage[utf8]{inputenc}
\usepackage{lmodern}
\usepackage[french]{babel}


\usepackage{graphicx}
\usepackage{tabularx}
\usepackage{minted}
\usepackage{eurosym}
\usepackage{enumitem}
\usepackage{amsmath}
\usepackage{tikz}


\usetikzlibrary{decorations.pathreplacing, backgrounds, calc}


\newcommand{\sql}[1] {\textsc{#1}}


\title{CC Term NSI - Arbre Binaire de Recherche}
\author{Babash - NSI Tani Malandi et Johanne Bakalara}
\date{Février 2025, Durée: 0h30min}


\usepackage[a4paper,top=2.3cm,bottom=2.8cm,left=2cm,right=2cm,headsep=0.6cm,footskip=1.2cm]{geometry}
\begin{document}
\runningfooter{}{\footnotesize Babash - NSI Tani Malandi et Johanne Bakalara \textperiodcentered{} CC BY-SA 4.0 \textperiodcentered{} page \thepage}{}
\firstpagefooter{}{\footnotesize Babash - NSI Tani Malandi et Johanne Bakalara \textperiodcentered{} CC BY-SA 4.0 \textperiodcentered{} page \thepage}{}

\maketitle
\thispagestyle{headandfoot}
\begin{center}\footnotesize\itshape D'après le sujet de NSI du baccalauréat 2023 (jour 2, Polynésie, code 23-NSIJ2PO1).\end{center}

\pagestyle{headandfoot}
\thispagestyle{headandfoot}
\firstpageheader{\makebox[\textwidth]{Nom : \dotfill\ Pr\'enom : \dotfill\ Classe : \makebox[3cm]{\dotfill}}}{}{}

\vsword{Note:}
\hsword{Note:}

\section*{Barème :}
\begin{center}
  \par\bigskip
  \gradetable[h]
\end{center}

\vfill
\underline{Arbre 1} \begin{center}
    \begin{tikzpicture}[node distance={1.5cm}, thick, main/.style = {draw, circle}] 
    \node[main] (20) {20}; 
    \node[main] (7) [below left of = 20, node distance={3.8cm}] {7}; 
    \node[main] (53) [below right of = 20, node distance={3.8cm}] {53}; 
    \node[main] (4) [below left of = 7] {4}; 
    \node[main] (15) [below right of = 7] {15}; 
    \node[main] (2) [below left of = 4] {2}; 
    \node[main] (11) [below left of =  15] {11}; 
    \node[main] (18) [below left of = 53] {18};
    \node[main] (80) [below  right of = 53] {80}; 
    \node[main] (83) [below right of=  80] {83}; 
    \draw (20) -- (7); 
    \draw (20) -- (53); 
    \draw (7) -- (4); 
    \draw (7) -- (15); 
    \draw (4) -- (2); 
    \draw (15) -- (11);
    \draw (53) -- (18);
    \draw (53) -- (80);
    \draw (80) -- (83);
    \end{tikzpicture} 
\end{center}

\vfill\newpage
\underline{Arbre 2} \begin{center}
    \begin{tikzpicture}[node distance={1.5cm}, thick, main/.style = {draw, circle}] 
    \node[main] (25) {25}; 
    \node[main] (10) [below left of = 25, node distance={3.5cm}] {10}; 
    \node[main] (40) [below right of = 25, node distance={3.5cm}] {40}; 
    \node[main] (5) [below left of = 10] {5}; 
    \node[main] (15) [below right of = 10] {15}; 
    \node[main] (30) [below left of = 40] {30}; 
    \node[main] (50) [below right of = 40] {50}; 
    \node[main] (7) [below left of = 15] {7}; 
    \node[main] (18) [below right of = 15] {18}; 
    \draw (25) -- (10); 
    \draw (25) -- (40); 
    \draw (10) -- (5); 
    \draw (10) -- (15); 
    \draw (40) -- (30); 
    \draw (40) -- (50);
    \draw (15) -- (7);
    \draw (15) -- (18);
    \end{tikzpicture} 
\end{center}


\vspace{2cm}
\underline{Arbre 3} \begin{center}
    \begin{tikzpicture}[node distance={1.5cm}, thick, main/.style = {draw, circle}] 
    \node[main] (7) {7}; 
    \node[main] (3) [below left of = 7, node distance={3.5cm}] {3}; 
    \node[main] (15) [below right of = 7, node distance={3.5cm}] {15}; 
    \node[main] (1) [below left of = 3] {1}; 
    \node[main] (5) [below right of = 3] {5}; 
    \node[main] (12) [below left of = 15] {12}; 
    \node[main] (20) [below right of = 15] {20}; 
    \node[main] (9) [below left of = 12] {9}; 
    \node[main] (13) [below right of = 12] {13}; 
    \node[main] (25) [below right of = 20] {25}; 
    \node[main] (8) [below left of = 9] {8}; 
    \draw (7) -- (3); 
    \draw (7) -- (15); 
    \draw (3) -- (1); 
    \draw (3) -- (5); 
    \draw (15) -- (12); 
    \draw (15) -- (20);
    \draw (12) -- (9);
    \draw (12) -- (13);
    \draw (20) -- (25);
    \draw (9) -- (8);
    \end{tikzpicture} 
\end{center}

\vspace{2cm}

\begin{questions}
\question[1] Quelle est la racine des arbres 1, 2 et 3 ?
\begin{description}
    \item[Arbre 1:]
    \item[Arbre 2:] 
    \item[Arbre 3:] 
\end{description}
\question[1] Quelles sont les feuilles des arbres 1, 2 et 3 ?
\begin{description}
    \item[Arbre 1:]
    \item[Arbre 2:] 
    \item[Arbre 3:] 
\end{description}
\question[1] Quelle est la taille des arbres 1, 2 et 3 ?
\begin{description}
    \item[Arbre 1:]
    \item[Arbre 2:] 
    \item[Arbre 3:] 
\end{description}
\question[1] Quelle est la profondeur du nœud 7 dans les arbres 1, 2 et 3 ?
\begin{description}
    \item[Arbre 1:]
    \item[Arbre 2:] 
    \item[Arbre 3:] 
\end{description}
\question[1] Quelle est la hauteur des arbres 1, 2 et 3 ?
\begin{description}
    \item[Arbre 1:]
    \item[Arbre 2:] 
    \item[Arbre 3:] 
\end{description}
\question[3] Dites pour chacun des arbres si c'est un Arbre Binaire de Recherche. Justifiez.
\begin{description}
    \item[Arbre 1:]
    \item[Arbre 2:] 
    \item[Arbre 3:] 
\end{description}
\question[3] Donner pour chacun des arbres le parcours en largeur.
\begin{description}
    \item[Arbre 1:]
    \item[Arbre 2:] 
    \item[Arbre 3:] 
\end{description}
\question[3] Donner pour chacun des arbres le parcours en profondeur préfixe.
\begin{description}
    \item[Arbre 1:]
    \item[Arbre 2:] 
    \item[Arbre 3:] 
\end{description}
\question[3] Donner pour chacun des arbres le parcours en profondeur infixe.
\begin{description}
    \item[Arbre 1:]
    \item[Arbre 2:] 
    \item[Arbre 3:] 
\end{description}
\question[3] Donner pour chacun des arbres le parcours en profondeur suffixe/postfixe.
\begin{description}
    \item[Arbre 1:]
    \item[Arbre 2:] 
    \item[Arbre 3:] 
\end{description}


\end{questions}
\end{document}